% Generated from validated Article Draft Result. Do not hand-edit; revise the structured draft instead.
\section{能力拡張と同時に増えた制御面――長文脈、behavior、interpretability}
\label{pkg:safety-control}
\noindent\textbf{長いcontextや強いtool能力が広がるほど、model behaviorをどう指定し、攻撃面をどう理解し、内部表現をどこまで観察できるかが重要になる。Many-shot jailbreaking、Model Spec、Anthropic interpretabilityを同じcontrol layerの異なる問題として読む。}\autocite{src-018fb3abd44d,src-4ff5ab7d45bc,src-d1406200c99e}

Many-shot jailbreakingは、context windowの拡大が便利さだけでなく新しいattack surfaceも生むことを示した。より多くのexampleを一度に与えられることは、few-shot promptingの強化と同時に、safety behaviorを崩すための入力空間も広げる。能力拡張と安全性は別々のrelease noteではなく、同じ設計変更の表裏として扱う必要がある。\autocite{src-018fb3abd44d,src-28ee6f0f00f4,src-e91e6f80eaac}


OpenAI Model Specは、modelに期待するbehaviorを明示的なspecificationとして表現する方向を示した。これはmodel内部を理解する方法ではなく、望ましい応答・拒否・優先順位を外部から定義するcontrol surfaceである。toolやagentへ接続されるほど、このbehavior layerは単なる会話品質ではなくsystem policyの一部になる。\autocite{src-d1406200c99e}


AnthropicのMapping the MindとGolden Gate Claudeは、mechanistic interpretabilityを通じてmodel内部のfeature表現を可視化・操作する研究方向を示した。behavior specificationが外側からルールを与えるのに対し、interpretabilityは内部表現を観察する。両者は補完的だが、内部featureを特定できることが一般的な安全制御を保証するわけではない。\autocite{src-4ff5ab7d45bc}


\begin{claimboundary}
jailbreak研究の結果は特定のthreat modelと実験条件に依存し、interpretability demonstrationも一般的なcontrol guaranteeではない。Model Specもbehaviorの意図を示す文書であり、すべてのmodel behaviorが常に仕様どおりになることを意味しない。本稿ではpolicy、attack evaluation、internal analysisを別レイヤとして扱う。\autocite{src-018fb3abd44d,src-4ff5ab7d45bc,src-5019b9256d83,src-d1406200c99e}
\end{claimboundary}

2024-H1では、能力競争の高速化に対してcontrol研究が後追いするだけではなく、長文脈、multimodality、agent化と並行して独立した設計層として立ち上がり始めた。modelを強くすることと、何をさせるか・なぜそう振る舞うかを理解することが、同じsystem engineeringの問題へ近づいた。\autocite{src-018fb3abd44d,src-28ee6f0f00f4,src-4ff5ab7d45bc,src-d1406200c99e,src-e91e6f80eaac}
